% !TeX root = ../../main.tex
\chapter{攻读博士学位期间的相关研究成果}
以下列出与本文直接相关的三项合作研究。会议与状态依提供的论文材料及作者确认记录整理，尚未确定的正式卷页与出版标识不作补写。
\begin{enumerate}
\item Jiaqi Zhai, Xuanhua Shi, Kaiyi Huang, Chencheng Ye, Weifang Hu, Bingsheng He, Hai Jin. AccelES: Accelerating Top-K SpMV for Embedding Similarity via Low-bit Pruning. IEEE International Symposium on High-Performance Computer Architecture (HPCA), 2025: 977--990.
\item Jiaqi Zhai, Xuanhua Shi, Wenju Zhao, Kaiyi Huang, Chencheng Ye, Shunsen Lv, Zhongtian Long, Bingsheng He, Hai Jin. ParetoES: Hardware-Accelerated Sparse Embedding Similarity via Pareto-Optimal Pruning. International Symposium on Computer Architecture (ISCA), 2026.
\item Jiaqi Zhai, Haozhi Han, Kaiyi Huang, Wenju Zhao, Yifeng Chen, Yunquan Zhang, Ting Cao, Yunxin Liu, Hai Jin, Xuanhua Shi, Kun Li. EvoCG: Self-Evolving Precision Control for Faster Conjugate Gradient. ACM SIGPLAN Symposium on Principles and Practice of Parallel Programming (PPoPP), 2027, 已录用。Jiaqi Zhai与Haozhi Han为共同第一作者。
\end{enumerate}

\chapter{研究成果与学位论文的关系}
\begin{table}[htbp]\centering
\caption{代表性研究成果与本文章节的对应关系}
\begin{tabularx}{\textwidth}{p{2cm}p{2.2cm}X}\toprule
研究成果&对应章节&在学位论文中的作用\\\midrule
AccelES&第3章&低比特分层计算、ReSparse、Ultra-CSR与Random-CSR、FPGA异构执行及实验\\
ParetoES&第4章&方向敏感聚类、候选簇选择、ACPE、分布式微排序及热点均衡\\
EvoCG&第5章&闭环精度控制、Tile-CSC、统一GPU内核、低开销运行时与数值验证\\
三项成果&第1、2、6章&重新组织共同研究问题、计算特征分析、方法关联、适用条件和未来方向\\\bottomrule
\end{tabularx}
\end{table}

三项成果均为合作研究。作者在具体算法设计、系统实现、实验分析和论文写作中的个人分工，应结合合作记录由作者和导师确认。本初稿不根据作者顺序推断独立完成比例，也不将合作团队完成的全部工作自动归属于本人。

第3章与第4章共享低比特稀疏数据流基础，本文在第4章明确区分继承技术与新增贡献。第5章研究具有不同数值语义的迭代过程，与前两章的关系是计算需求驱动方法的拓展，不是同一算法或加速器的连续版本。

\chapter{主要符号及中英文缩写}
\begin{longtable}{p{2.7cm}p{5.3cm}p{5.3cm}}
\toprule 缩写或符号&英文或数学含义&中文含义\\\midrule\endhead
SpMV&Sparse Matrix--Vector Multiplication&稀疏矩阵--向量乘法\\
Top-K&Top-K selection&最高分的K个结果选择\\
CG&Conjugate Gradient&共轭梯度法\\
PCG&Preconditioned Conjugate Gradient&预条件共轭梯度法\\
SPD&Symmetric Positive Definite&对称正定\\
CSR&Compressed Sparse Row&压缩稀疏行格式\\
CSC&Compressed Sparse Column&压缩稀疏列格式\\
COO&Coordinate format&坐标格式\\
NNZ&Number of nonzeros&非零元素数量\\
HBM&High Bandwidth Memory&高带宽存储器\\
FPGA&Field-Programmable Gate Array&现场可编程门阵列\\
GPU&Graphics Processing Unit&图形处理器\\
INT6&6-bit signed integer&6位有符号整数\\
FP32 / FP64&32/64-bit floating point&单精度／双精度浮点\\
BF16&Brain Floating Point 16&16位浮点格式\\
ACPE&Adaptive Cluster Probing Engine&自适应簇探测引擎\\
\hbalance{}&Hierarchical Hotspot-Balancing&分层热点负载均衡\\
PPO&Proximal Policy Optimization&近端策略优化\\
QPS&Queries Per Second&每秒查询数\\
CV&Coefficient of Variation&变异系数\\
$m,n,z$&Rows, columns, nonzeros&行数、列数、非零元素数\\
$K,L$&Result size, candidate budget&输出规模、预筛候选预算\\
$R,C$&Tile height and width&Tile行数与列数\\
$r_k,p_k$&Residual and search direction&第k轮残差与搜索方向\\
$K^*$&Action reuse interval&精度动作复用间隔\\\bottomrule
\end{longtable}
